\section{Linear Search: an Iteration Constraint Abstracted and Instantiated}\label{sec:lsearch}

\subsection{The program}
\begin{lstlisting}
method lsearch(a: array<int, int>, n: int, x: int) returns (i: int)
  requires 0 <= n
  ensures 0 <= i <= n
  ensures i < n ==> a[i] == x
  ensures forall k: int :: 0 <= k < i ==> a[k] != x
{
  i := 0;
  while i < n && a[i] != x { i := i + 1; }
}
\end{lstlisting}
The first index of $x$ among $a[0..n)$, or $n$. The loop test is data-dependent, so the number of
iterations is not determined by $n$; the trace equation system records the test as the iteration
constraint and its negation as the exit.

\subsection{The trace equation system}
\begin{align}
  K &= \cnt(0, 1, n') \qquad i = \cur(I) \qquad I = \tr{0}{\bl{K}}{i, k \mapsto i + 1} \label{eq:ls0}\\
  &\al{\bl{K}, \bl{I}}{k, i \mapsto i < n \wedge a[i] \ne x} \label{eq:ls1}\\
  &\neg(i < n \wedge a[i] \ne x) \label{eq:ls2}\\
  &\Longrightarrow\ 0 \le i \le n\ \wedge\ (i < n \Rightarrow a[i] = x)\ \wedge\ (\forall k.\ 0 \le k < i \Rightarrow a[k] \ne x) \nonumber
\end{align}
With \texttt{-split}, the conclusion becomes four goals, one per conjunct, the last two with their premises
as hypotheses: $i < n \Rightarrow a[i] = x$, and $0 \le k < i \wedge a[k] = x \Rightarrow \mathsf{false}$.
The simplifier's own rules do the following: $I$ is a counter by \lemma{counter\_def}, so
$i = \intf(n')$; the constraint \eqref{eq:ls1} is distributed into
$\al{\bl{\cnt(0, 1, n')}}{k \mapsto k < n}$, which \lemma{invs\_counter\_upper} resolves to
$n' = 0 \vee \intf(n') \le n$ (hence $0 \le i \le n$, the first two goals), and
$\al{\bl{\cnt(0, 1, n')}}{k \mapsto a[k] \ne x}$, the search constraint proper; the exit
\eqref{eq:ls2} is $n \le i \vee a[i] = x$, which with $i < n$ gives the third goal.

\subsection{The lemmas}
The last goal needs the search constraint at a witness: for $0 \le k < i$, the constraint says
$a[k] \ne x$. The paper's lemma database has no instantiation of iteration constraints at positions
(\texttt{TODO.md}: no instantiation at a witness); the two lemmas of \texttt{lsearch.dfy}, since merged
into \texttt{teq.dfy}, abstract the constraint by an atom and instantiate the atom:
\begin{align*}
  &\lemma{not\_occurs\_selects\_def}:\\
  &\quad \al{\bl{\cnt(i, 1, n)}}{k \mapsto a[k] \ne x}\ \Longleftrightarrow\ \neg\mathsf{occurs}(x, \mathsf{selects}(a, \bl{\cnt(i, 1, n)}))\\[4pt]
  &\lemma{select\_not\_occurs\_selects}:\\
  &\quad \neg\mathsf{occurs}(x, \mathsf{selects}(a, \bl{\cnt(i, 1, n)}))\ \wedge\ i \le k < i + \intf(n)
    \ \Longrightarrow\ (a[k] = x) \Longleftrightarrow \mathsf{false}
\end{align*}
The first matches the constraint with $i := 0$, $n := n'$ (the counter $K$ is followed in the dependency
position) and replaces it by the atom "$x$ does not occur among the values of $a$ at the positions
$0, \ldots, n' - 1$". The second is a rule on the atom $a[k] = x$ of the last goal: its side conditions
are the fact just produced, matched with $i := 0$, $n := n'$, and the bounds $0 \le k < \intf(n')$,
decided from $0 \le k < i$ and $i = \intf(n')$. The hypothesis $a[k] = x$ becomes $\mathsf{false}$ and
the goal closes.

This is the pattern of \lemma{not\_occurs\_prefix\_def} and \lemma{earlier\_not\_occurs\_prefix} of
\texttt{teq.dfy} for reads $\mathsf{earlier}(k, xs)$ of a sequence, here for reads of an array: an
iteration constraint becomes a named fact about the prefix visited, and a second lemma uses the fact at
any position in the prefix. \texttt{TODO.md} notes that a side condition that is the iteration constraint
itself would make the rule cycle (finding 19 of \texttt{NOTES.md}), which is why the abstraction by an
atom is not merely convenient.

The file's header prescribes \texttt{-teq -split -simplify}, since without the split the disjunction of
the two exits, $n \le i \vee a[i] = x$, had to be split by the simplifier against the postcondition; with
the current version the goal closes without \texttt{-split} as well.
